\documentclass[english, parskip=full]{scrartcl}
\usepackage[utf8]{inputenc}  % Kodierung der Datei
\usepackage[english]{babel}  % Sprache auf Deutsch 
\usepackage{color}
\usepackage{graphicx, gensymb}
\usepackage{amsfonts, cancel, amssymb}
\usepackage{amsmath, amsthm, array, esint}
\usepackage{booktabs, multirow}  % schönere Tabellen
\usepackage{caption}  % Anpassung der Captions
\usepackage{scrlayer-scrpage}    % Manipulation des Seitenstils
\pagestyle{scrheadings}  % Stil für die Seite setzen
\automark{section}  % Abschnittsnamen als Seitenbeschriftung 
\setheadsepline{.5pt}    % Kopzeile durch Linie abtrennen
\usepackage[hidelinks]{hyperref}  % Links und weitere PDF-Features
\usepackage{typearea, tikz, pgffor, verbatim}
\usetikzlibrary{calc, arrows.meta,arrows, graphs, quotes, positioning}
\usepackage[cache=false, outputdir=build]{minted}
\usemintedstyle{vs}
\usepackage{placeins} % provides \FloatBarrier

\definecolor{bg}{rgb}{0.9,0.9,0.9}
\newcommand\numberthis{\addtocounter{equation}{1}\tag{\theequation}}
\usepackage{svg}
\usepackage{siunitx}
\usepackage{physics}
\usepackage{tensor}
\renewcommand{\bra}{}
\newcommand{\kom}{}
\newcommand{\brak}{}
\renewcommand{\braket}{}
\newcommand{\intx}{}
\newcommand{\intr}{}
\newcommand{\kla}{}
\newcommand{\klam}{}
\newcommand{\klammer}{}
\newcommand{\oper}{}
\newcommand{\tecxt}{}
\newcommand{\Bra}{}
\newcommand{\Ket}{}
\renewcommand{\i}{\text{i}}
\newcommand{\e}{\text{e}}
\newcommand*{\Fourier}{\textsc{Fourier}\ }
\newcommand*{\F}{\mathcal{F}}
\newcommand*{\sinc}{\text{sinc\,}}
\newcommand{\conv}{\mathop{\scalebox{3.5}{\raisebox{-0.38ex}{\makebox[0.5\width][c]{$\ast$}}}}\limits}

\newcommand*{\titel}{Bose-Einstein Condensation in 2D}
\newcommand*{\autor}{Joachim Schwardt}

\hypersetup{pdfauthor={\autor}, pdftitle={\titel}}

%\titlehead{titlehead}
\subject{Theory of Superconductivity}
\title{\titel}
\author{\autor}
\date{\today}

\begin{document}

%\begin{titlepage}
\maketitle  % Titel setzen
%\tableofcontents  % Inhaltsverzeichnis setzen
%\end{titlepage}

%\section{London penetration depth}
%In gaussian units we have and assuming $\rho \equiv 0$ the Maxwell and London equations in real and Fourier space read
%\begin{align}
%\nabla\cdot \textbf{E} &= 0 && \rightarrow & \i\textbf{k} \cdot \textbf{E} &= 0,\\
%\nabla\cdot \textbf{B} &= 0 && \rightarrow & \i\textbf{k} \cdot \textbf{B} &= 0,\\
%\nabla\times \textbf{E} &= - \frac{1}{c} \partial_t \textbf{B} && \rightarrow & \textbf{k} \times \textbf{E} &= \frac{1}{c} \omega \textbf{B}, \label{eqn:maxwell 3}\\
%\nabla\times \textbf{B} &= \frac{4\pi}{c}\textbf{j} + \frac{1}{c} \partial_t \textbf{E} && \rightarrow & \textbf{k} \times \textbf{B} &= \frac{4\pi}{\i c}\textbf{j} - \frac{\omega}{c} \textbf{E},\\
%\partial_t \textbf{j} &= \frac{c^2}{4\pi\lambda_L^2} \textbf{E} && \rightarrow & \i\omega \textbf{j} &= \frac{c^2}{4\pi\lambda_L^2} \textbf{E}, \\
%\nabla\times \textbf{j} &= -\frac{c}{4\pi\lambda_L^2} \textbf{B} && \rightarrow & \i\textbf{k} \times \textbf{j} &= -\frac{c}{4\pi\lambda_L^2} \textbf{B}.
%\end{align}
%The solutions have to be transverse waves since both $\textbf{k} \cdot \textbf{E} = 0$ and $\textbf{k} \cdot \textbf{B} = 0$.
%\autoref{eqn:maxwell 3} can be written as $\textbf{B} = \frac{c}{\omega} \textbf{k} \times \textbf{E}$.
%Using this, we may eliminate $\textbf{B}$ from the set of equations
%\begin{align}
%\frac{4\pi}{\i c} \textbf{j} - \frac{\omega}{c} \textbf{E} &= \textbf{k} \times \kla\left( \frac{c}{\omega} \textbf{k} \times \textbf{E} \right) = \frac{c}{\omega} \textbf{k} \kla\left( \textbf{k} \cdot \textbf{E} \right) - \frac{ck^2}{\omega} \textbf{E} =  - \frac{ck^2}{\omega} \textbf{E},\\
%\i \textbf{j} &= -\frac{c}{4\pi\lambda_L^2} \frac{c}{\omega} \textbf{E}.
%\end{align}
%We find two equations for the supercurrent $\textbf{j}$
%\begin{align}
%\textbf{j} &= \frac{\i}{4\pi} \left( \omega - \frac{c^2 k^2}{\omega} \right) \textbf{E} ,\\
%\textbf{j} &= \frac{\i c^2}{4\pi\omega\lambda_L^2} \textbf{E}.
%\end{align}
%Combining these leads to an equation for the dispersion relation
%\begin{align}
%0 &= \omega - \frac{c^2k^2}{\omega} - \frac{c^2}{\omega\lambda_L^2} && \Rightarrow & \omega &= c\sqrt{k^2 + \frac{1}{\lambda_L^2}} \equiv \sqrt{c^2k^2 + \omega_p^2}
%\end{align}
%where $\omega_p := \frac{c}{\lambda_L}$ is the plasma frequency.
%\autoref{fig:long dispersion} shows a plot of the dispersion relation.
%\begin{figure}[!ht]
%\centering
%\includegraphics{london_dispersion.png}
%\caption{Plot of the dispersion relation for $\lambda_L = \SI{200}{nm}$.}
%\label{fig:long dispersion}
%\end{figure}
%\section{Boltzmann equation in relaxation time approximation}
%In the absence of an applied force the Boltzmann equation in the relaxation time approximation $\mathcal{S}[\rho] = \frac{\rho - \rho_0}{\tau}$ reads
%\begin{align}
%\kla\left( \partial_t + \frac{\hbar\textbf{k}}{m} \cdot\nabla \right) \rho &= - \frac{\rho - \rho_0}{\tau}
%\end{align}
%where $\rho_0(k) = \frac{1}{\e^{\beta \epsilon_\textbf{k}} + 1}$ is the Fermi-Dirac distribution function.
%Since both $\partial_t \rho_0 = 0$ and $\nabla\rho_0 = 0$ we indeed find that $\rho = \rho_0$ is a (trivial) solution to this simplified Boltzmann equation.
%Let $\eta(\textbf{r},\textbf{k},t) := \rho (\textbf{r},\textbf{k},t) - \rho_0(\textbf{k})$ be the deviation from the equilibrium distribution.
%Then the equation for $\eta$ reads
%\begin{align}
%\kla\left( \partial_t + \frac{\hbar \textbf{k}}{m} \cdot \nabla \right) \eta &= -\frac{\eta}{\tau}.
%\end{align}
%%Applying a spatial Fourier transform results in
%%\begin{align}
%%\kla\left( \partial_t + \frac{\hbar \textbf{k}}{m} \cdot \i \textbf{q} \right) \eta(\textbf{q},\textbf{k},t) &= -\frac{\eta(\textbf{q},\textbf{k},t) }{\tau}.
%%\end{align}
%%Splitting of the exponential dampening term by $\eta =: \phi (\textbf{q},\textbf{k}) \e^{-t/\tau}$ yields
%%\begin{align}
%%\frac{\hbar}{m\tau} \textbf{k} \cdot \i\textbf{q} \phi &= 0.
%%\end{align}
%Splitting of the exponential dampening term by $\eta =: \phi (\textbf{r},\textbf{k}) \e^{-t/\tau}$ yields
%\begin{align}
%\frac{\hbar}{m\tau} \textbf{k} \cdot \nabla \phi &= 0.
%\end{align}
%A simple solution of this equation would be given by an arbitrary, but spatially homogeneous $\phi = \phi(\textbf{k})$.
%Regardless, the deviations show exponential dampening with an expected lifetime of $\tau$.

%\section{Bose-Einstein condensation in two dimensions}
This problem is best solved in the grand-canonical ensemble where
\begin{align}
\hat{\rho} &= \frac{1}{Z} \e^{-\beta \hat{H} + \beta\mu \hat{N}} = \frac{1}{Z} \prod_\textbf{k} \e^{-\beta \epsilon_\textbf{k} \hat{n}_\textbf{k} + \beta \mu \hat{n}_\textbf{k}}.
\end{align}
Since the trace of this density operator is unity the partition function is given by
\begin{align}
Z &= \tecxt\text{tr} \prod_\textbf{k} \e^{-\beta \epsilon_\textbf{k} n_\textbf{k} + \beta \mu n_\textbf{k}} = \sum_{\{n_\textbf{k}\}} \langle \{n_\textbf{k}\} | \prod_\textbf{q} \e^{-\beta (\epsilon_\textbf{q} - \mu) \hat{n}_\textbf{q}} | \{n_\textbf{k} \} \rangle = \\
&= \sum_{\{n_\textbf{k}\}} \prod_\textbf{q} \e^{-\beta (\epsilon_\textbf{q} - \mu) n_\textbf{q}} = \prod_\textbf{k} \sum_{n_\textbf{k} \ge 0} \e^{-\beta (\epsilon_\textbf{k} - \mu) n_\textbf{k}} = \prod_\textbf{k} \frac{1}{1 - \e^{-\beta (\epsilon_\textbf{k} - \mu)}}.
\end{align}
With the fugacity $y := \e^{\beta\mu}$ the logarithm of this partition function becomes
\begin{align}
\log Z &= \sum_\textbf{k} \log \kla\left( 1 - y\e^{-\beta\epsilon_\textbf{k}} \right).
\end{align}
The expected number of particles is given by
\begin{align}
N &= \tecxt\text{tr} (\hat{N} \hat{\rho}) = \frac{1}{Z} \frac{\partial_\mu }{\beta} Z = \frac{1}{\beta} \partial_\mu \log Z = y\partial_y \log Z = \sum_\textbf{k} \frac{y\e^{-\beta\epsilon_\textbf{k}}}{1 - y\e^{-\beta\epsilon_\textbf{k}}}.
\end{align}
The density of states is defined by
\begin{align}
n(\epsilon) &:= \frac{1}{A}\sum_\textbf{k} \delta(\epsilon - \epsilon_\textbf{k}).
\end{align}
This allows us to rewrite expressions of the following form in terms of an integral
\begin{align}
%\sum_\textbf{k} \cdot &\sim \frac{A}{(2\pi)^2} \intr\int_{\mathbb{R}^{2}} \cdot\ \,\mathrm{d}^{2}k = \frac{A}{(2\pi)^2} \intx\int_{0}^{\infty} \cdot  \  n(\epsilon) \, \mathrm{d}\epsilon
\sum_\textbf{k} f(\epsilon_\textbf{k}) &= \sum_\textbf{k} \intx\int_{0}^{\infty} f(\epsilon) \delta (\epsilon - \epsilon_\textbf{k}) \, \mathrm{d}\epsilon = A\int_{0}^{\infty} f(\epsilon) n(\epsilon)\, \mathrm{d}\epsilon.
\end{align}
In our case the density of states may be approximated by
\begin{align}
n(\epsilon) &\simeq \frac{1}{(2\pi)^2} \intr\int_{\mathbb{R}^{2}} \delta \kla\left( \epsilon - \frac{\hbar^2\textbf{k}^2}{2m} \right) \, \mathrm{d}^{2}k = \frac{1}{2\pi} \intx\int_{0}^{\infty} \delta \kla\left( \epsilon - \frac{\hbar^2k^2}{2m} \right) \, \frac{1}{2}\mathrm{d}(k^2) = \frac{m}{2\pi \hbar^2}.
\end{align}
Thus we find
\begin{align}
\log Z &= \frac{mA}{2\pi\hbar^2} \intx\int_{0}^{\infty} \log \kla\left( 1 - y\e^{-\beta\epsilon} \right) \, \mathrm{d}\epsilon,\\
N &= \frac{mA}{2\pi\hbar^2} \intx\int_{0}^{\infty} \frac{1}{\frac{1}{y} \e^{\beta\epsilon} - 1}\, \mathrm{d}\epsilon = \frac{mA}{2\pi\hbar^2\beta} \intx\int_{0}^{\infty} \frac{y\e^{-x}}{1 - y\e^{-x}}\, \mathrm{d}x = \\
&= \frac{mA}{2\pi\hbar^2\beta} \sum_{n\ge 1} \intx\int_{0}^{\infty} \kla\left( y\e^{-x} \right)^n\, \mathrm{d}x = \\
&=\frac{mA}{2\pi\hbar^2\beta} \sum_{n\ge 1} y^n \frac{1}{n} = \frac{mA}{2\pi\hbar^2\beta} \log \kla\left( \frac{1}{1 - y} \right).
\end{align}
The fugacity can now be expressed as a function of the particle number
\begin{align}
y &= 1 - \e^{-\frac{2\pi\hbar^2\beta}{m} \frac{N}{A}}
\end{align}
and similarly the partition function becomes
\begin{align}
\log Z &= \frac{mA}{2\pi\hbar^2} \intx\int_{0}^{y} \log\kla\left( 1 - u \right) \, \frac{\mathrm{d}u}{\beta u} = \\
&= -\frac{mA}{2\pi\hbar^2\beta} \tecxt\text{Li}_2(y) = -\frac{mA}{2\pi\hbar^2\beta} \tecxt\text{Li}_2 \kla\left( 1 - \e^{-\frac{2\pi\hbar^2}{mA} \beta N} \right).
\end{align}
Note that the ground state occupation number is $N_0 = \frac{y}{1 - y}$.
For a phase transition this would have to diverge, requiring $y\rightarrow 1$ and therefore $\beta\mu \rightarrow 0$.
Thus the chemical potential would have to scale as $\mu \sim T^{1+\alpha}$ for some $\alpha > 0$.
Comparing to our equation for the particle density, this would be equivalent to
\begin{align}
\frac{N}{A} \frac{2\pi\hbar^2}{mk_\tecxt\text{B}} &= -T \log \kla\left( 1 - \e^{\kappa T^{1 + \alpha}} \right) \\
&\rightarrow \lim \frac{\frac{1}{1 - \e^{\kappa T^{1+\alpha}}} \kla\left( -\kappa (1+\alpha) T^\alpha \e^{\kappa T^{1+\alpha}} \right)}{\frac{1}{T^2}} = \lim \frac{-\kappa (1+\alpha) T^{2+\alpha}}{1 - \e^{\kappa T^{1+\alpha}}} = 0.
\end{align}
However, since we are dealing with bosons we expect all particles to be in the ground state for $T=0$.
There is no phase transition, because as seen in the limit above, the particle density of the $N_{>0}$-states non-zero for any finite temperature.
%Note that $\mu \le 0$ and therefore $y = \e^{\beta\mu} \rightarrow 0$ for $T \rightarrow 0$ (or $\beta\rightarrow \infty$ respectively).
%Thus
%\begin{align}
%\frac{N}{A} &= \frac{-m}{2\pi\hbar^2\beta} \log (1 - y) \sim \frac{m}{2\pi\hbar^2\beta} y = \frac{m}{2\pi\hbar^2} \frac{\e^{\beta\mu}}{\beta} \le \frac{m}{2\pi\hbar^2\beta} \rightarrow 0.
%\end{align}
%\begin{align}
%\mu &= \frac{1}{\beta} \log y = \frac{1}{\beta} \log \kla\left( 1 - \e^{-\frac{2\pi\hbar^2\beta}{m} \frac{N}{A}} \right) \le 0,\\
%\frac{N}{A} &\le \frac{m}{2\pi\hbar^2} \frac{y}{\beta} \le \frac{m}{2\pi\hbar^2} \frac{1}{\beta} \rightarrow 0.
%\end{align}
%TODO: Show that there is no phase transition (this is more difficult than it should be, something weird about this approach) and compute pressure (should be easy).
The grand-canonical potential is $\Omega = -\frac{1}{\beta} \log Z$.
The pressure is then given by
\begin{align}
p &= -\frac{\partial \Omega}{\partial A} = -\frac{\Omega}{A} + \frac{1}{\beta} \frac{mA}{2\pi\hbar^2\beta} \frac{-\log (1-y)}{y} \frac{\partial y}{\partial A} = \\
&= -\frac{\Omega}{A} + \frac{1}{\beta} \frac{mA}{2\pi\hbar^2\beta} \frac{-\log (1-y)}{y} \frac{-2\pi\hbar^2\beta N}{mA^2} (1 - y) = \\
&= \frac{m}{2\pi\hbar^2\beta^2} \tecxt\text{Li}_2(y) + \frac{1}{\beta A} \frac{\log (1-y)}{y}(1-y).
\end{align}
The low-temperature limit is equivalent to $y\rightarrow 1$ for which the dilogarithm behaves as $\tecxt\text{Li}_2(y) \sim \text{Li}_2(1) - (1-y) \klam\left[ 1 - \log(1 - y) \right]$.
Thus the pressure becomes
\begin{align}
p &\sim \frac{m}{2\pi\hbar^2\beta^2} \klam\left[ \text{Li}_2(1) - (1-y) \right] + (1-y)\log(1-y) \klam\left[ \frac{m}{2\pi\hbar^2\beta^2} + \frac{1}{\beta A y} \right]\\
&\rightarrow \frac{m}{2\pi\hbar^2\beta^2} \frac{\pi^2}{6} = \frac{\pi m}{12 \hbar^2}(k_\tecxt\text{B}T)^2.
\end{align}
In a classical ideal gas $p\sim T$ and in a 3D condensate $p\sim T^{5/2}$, so pressure in 2D drops faster than in a classical gas but slower than in 3D.

\subsubsection*{Problems with $T\rightarrow 0$}
The script and problem sheet both claim that ``splitting off the ground state would be not just unnecessary but incorrect for a two-dimensional gas''. 
But why would this be case?
Either the ground state shows macroscopic occupation (then we do have to split it off) or it does not.
In case of the latter we certainly have
\begin{align}
N &= (\textbf{k} = 0)-\tecxt\text{term} + A \int \dots
\end{align}
where, since we assumed the first term to not be macroscopically occupied,
\begin{align}
\frac{N}{A} &= \frac{(\textbf{k} = 0)-\tecxt\text{term}}{A} + \int \dots \overset{\tecxt\text{t.d. limit: } N,A \rightarrow \infty}{\longrightarrow} 0 + \int \dots
\end{align}
so the ``wrong'' term just disappears again.
Another line of thought is that, unless one is dealing with divergences, splitting off zero-measure sets from the realm of integration does not change the result of an integral.
Of course my approach must be flawed in some way, but I did not find the mistake.
%\\
%Regardless of this issue (which is presumably due to a major misunderstanding on my part of course\dots), I also fail to see how the situation for the $\textbf{k} = 0$-term is any different from the 3D case since $\epsilon_0 = 0$ all the same and therefore
%\begin{align}
%N_0 &= \frac{y}{1 - y} 
%\end{align}
%may still diverge for $T\rightarrow 0$ if $y \rightarrow 1$.
%However, this would require $\beta\mu \rightarrow 0$ and thus $\mu \sim T^{1+\alpha}$ for some $\alpha > 0$.
%\begin{align}
%\frac{N}{A} &= \int_{\mathbb{R}^{2}} \frac{1}{\e^{\beta (\epsilon_\textbf{k} - \mu)} - 1} \, \frac{\mathrm{d}^{2}k}{(2\pi)^2} = \frac{1}{(2\pi)^2} \intx\int_{0}^{\infty} \frac{1}{\e^{\beta (\epsilon - \mu)} - 1} \, n(\epsilon)\mathrm{d}\epsilon = \\
%&= \frac{m}{(2\pi)^3\hbar^2\beta} \intx\int_{-\beta\mu}^{\infty} \frac{1}{\e^x - 1} \, \mathrm{d}x = \\
%&= \frac{m}{(2\pi)^3\hbar^2\beta} \zeta(2) + \frac{m}{(2\pi)^3\hbar^2\beta} \intx\int_{-\beta\mu }^{0} \frac{1}{\e^x - 1}\, \mathrm{d}x \sim \\
%&\sim \frac{mk_\tecxt\text{B} }{48\pi\hbar^2}T + \frac{m}{(2\pi)^3\hbar^2\beta} \intx\int_{-\beta\mu}^{0} \frac{1}{x}\, \mathrm{d}x \sim \dots + \dots \frac{\log}{•}
%\end{align}


%\begin{thebibliography}{9}
%\bibitem{example} description, \url{https://www.exmapleurl}, day.\,month~2021.
%\end{thebibliography}

\end{document}